% Обзор литературы 0001.
% Сборка: xelatex review.tex && bibtex review && xelatex review.tex && xelatex review.tex
% PDF кладётся в out/, в репозиторий не коммитится.
\documentclass[11pt,a4paper]{article}

\usepackage{iftex}
\ifPDFTeX
  \usepackage[T2A]{fontenc}
  \usepackage[utf8]{inputenc}
  \usepackage[russian]{babel}
\else
  \usepackage{fontspec}
  \usepackage[russian]{babel}
  \IfFontExistsTF{Times New Roman}
    {\setmainfont{Times New Roman}}
    {\IfFontExistsTF{DejaVu Serif}
      {\setmainfont{DejaVu Serif}}
      {\IfFontExistsTF{Liberation Serif}
        {\setmainfont{Liberation Serif}}
        {}}}
\fi

\usepackage[margin=2.4cm]{geometry}
\usepackage{amsmath}
\usepackage{booktabs}
\usepackage{longtable}
\usepackage{array}
\usepackage{enumitem}
\usepackage{hyperref}
\usepackage{microtype}

\hypersetup{
  colorlinks=true,
  linkcolor=black,
  citecolor=black,
  urlcolor=blue
}

\setlist[itemize]{leftmargin=1.4em,itemsep=0.2em,topsep=0.2em}
\setlist[enumerate]{leftmargin=1.6em,itemsep=0.2em,topsep=0.2em}

\title{Априорная оценка сложности задачи по её условию:\\
обзор литературы}
\author{Обзор 0001}
\date{8 ноября 2026}

\begin{document}
\maketitle

\begin{abstract}
\noindent
Обзор отвечает на вопрос, как до начала исполнения оценить сложность задачи
для LLM-агента: вероятность успеха, расход токенов и выполнимость. Включена
41 работа: прямая априорная оценка для агентских задач, психометрика и IRT,
предсказание сложности задач по программированию, маршрутизация запросов,
распределение бюджета и классическое предсказание времени работы алгоритма.
Вывод: по отдельности каждая величина предсказуема, но совместной оценки
«успех + стоимость» по условию задачи для кодового агента при общем бюджете в
литературе нет. Сильнее всего работают структурные признаки задачи, а не
лингвистические; прямая оценка сложности сильной LLM без калибровки
ненадёжна.
\end{abstract}

\section{Вопрос и зачем}

\subsection{Вопрос}

Как по условию задачи, доступному до начала исполнения, оценить её сложность
для LLM-агента: вероятность успеха, потребление токенов и выполнимость?

Слово «сложность» здесь распадается на три величины, которые литература
изучает раздельно:

\begin{enumerate}
  \item $p(\text{успех} \mid \text{задача}, \text{агент})$ — вероятность, что
        конкретная пара «модель + scaffold» решит задачу;
  \item $C$ — стоимость решения в токенах, шагах или секундах;
  \item $F$ — выполнимость: решаема ли задача вообще данным агентом.
\end{enumerate}

Третья величина не сводится ко второй: задача может быть дорогой, но
решаемой, а может быть нерешаемой при любой трате бюджета. Обзор показывает,
что и методы для $p$, $C$ и $F$ разные.

\subsection{Зачем}

Харнесс получает 100 задач, общий бюджет токенов и баллы за задачи. Баллы
расставлены случайно и не связаны со сложностью, 15--25\% задач невыполнимы.
Сейчас харнесс идёт в лексикографическом порядке каталогов, не оценивает
расход и не отличает невыполнимые задачи. Ответ на вопрос определяет, можно
ли заменить фиксированный порядок выбором задач по ожидаемой выгоде на токен
и стоит ли писать спецификацию изменения
\texttt{04\_budget\_and\_selection.md}.

Обзор не проектирует изменение. Он фиксирует, что уже проверено, какие
метрики использовались и где остаётся пробел.

\section{Метод}

\subsection{Базы и запросы}

Поиск выполнен 8 ноября 2026 года в двух базах: arXiv API (поле
\texttt{abs:}) и OpenAlex (фильтр \texttt{title\_and\_abstract.search}).
Просматривались верхние 8--25 записей на запрос, сортировка по релевантности.
Язык — английский, годы — 2012--2026. Более ранние работы не включались,
кроме метода предсказания времени работы алгоритма, на который прямо
опираются современные работы.

Semantic Scholar и Google Scholar были недоступны (HTTP 403), поэтому
использованы две другие базы. Точные запросы, даты и числа приведены в
\texttt{protocol.md}.

\subsection{Критерии отбора}

Включались работы, которые:
\begin{itemize}
  \item предсказывают сложность, стоимость, успех или выполнимость задачи до
        исполнения либо дают переносимый метод (IRT, предсказание времени
        работы алгоритма);
  \item имеют количественную оценку качества или стратегии;
  \item относятся к текстовым, кодовым, агентским или плановым задачам;
        образовательная психометрика допускалась, если сложность
        предсказывается по тексту задания.
\end{itemize}

Исключались работы, где сложность измеряется только постфактум, где оценка
требует недоступных до исполнения сигналов (несколько прогонов, скрытые
состояния, сенсоры) и где нет доступа к тексту.

\subsection{Отбор}

Из 339 идентифицированных записей после удаления 8 внутренних дубликатов
просмотрена 331. На скрининге исключена 271 запись, для разбора отобрано 60,
после разбора исключено 19. В обзор включена 41 работа. Схема и причины
исключения — в \texttt{sources.md}.

\subsection{Извлечение и оценка качества}

Для каждой работы выписаны цель, признаки, метод, данные, метрика, результат и
ограничение. Надёжность оценивалась по рецензируемости, размеру данных,
наличию отложенной выборки, сравнению с baseline и обработке cold start.
Мета-анализ не проводился: метрики и данные несопоставимы.

\section{Что известно}

\subsection{Прямая априорная оценка для агентских задач}

Девять работ изучают именно ту постановку, которая нужна харнессу: предсказать
исход задачи по её описанию до запуска агента.

\textbf{Сложность предсказуема по статическим признакам.}
Al-Haque и Johnson~\cite{alhaque2026difficult} на крупнейшем открытом корпусе
траекторий кодового агента CoderForge-Preview получили AUC~$=0{,}863$ для
предсказания успеха issue-resolution задачи. Главные факторы — фрагментация
патча и масштаб репозитория; лингвистические признаки промпта выходят вперёд
только в средней полосе сложности. Это прямой аргумент, что условие задачи
несёт сигнал о сложности, причём структурный, а не только текстовый.

\textbf{Оценка должна быть относительной.}
Ge и соавторы~\cite{ge2026psychometrics} расширяют IRT признаками задачи
(issue statement, репозиторий, решение, тесты) и раскладывают способность
агента на компоненты LLM и scaffold. Модель предсказывает успех на невиданных
задачах и невиданных комбинациях «модель--scaffold». Отсюда следует важное для
нас ограничение: сложность задачи не является свойством только текста, она
определяется парой «задача--агент». Оценку, снятую на одной модели, нельзя
переносить на другую без калибровки.

\textbf{Структура условия информативнее текста.}
CIRCE~\cite{hong2026circe} извлекает из условия граф ограничений и предсказывает
правильность ответа LLM: AUROC растёт с $0{,}691$ до $0{,}727$, NLL падает с
$0{,}465$ до $0{,}423$ на семи бенчмарках. TaskSense~\cite{yin2025tasksense}
строит «когнитивную цепочку» GUI-задачи и объясняет $R^2 = 0{,}46$ дисперсии
времени выполнения. Оба результата получены без исполнения задачи.

\textbf{Стоимость можно оценить теоретически.}
Lutati~\cite{lutati2025capability} формулирует Agent Capability Problem:
агенту нужно $I_{\text{total}}$ бит, чтобы найти решение, и он получает
$I_{\text{step}}$ бит за действие ценой $C_{\text{step}}$. Эффективная
стоимость $C_{\text{eff}} = (I_{\text{total}}/I_{\text{step}}) \cdot
C_{\text{step}}$ нижне ограничивает ожидаемые затраты. Это единственная
найденная работа, которая даёт априорную оценку стоимости, а не только
вероятности успеха, но она требует оценить информационную структуру задачи.

\textbf{Метрики обманчивы.}
Krsteski и Meyer~\cite{krsteski2026difficulty} на 17 агентских бенчмарках
показывают, что AUC может маскировать плохие оценки сложности; полезным
сигналом оказывается энтропия токенов при чтении условия. Ещё резче это
формулируют EsfandyariDoulabi и соавторы~\cite{esfandiaridoulabi2026disentangling}:
более 99{,}93\% пар «успех--провал» в pooled-оценках относятся к разным
задачам, а не к разным прогонам одной задачи. Поэтому AUROC $0{,}85$--$0{,}94$
у предикторов провала в основном измеряет сложность задачи, а не приближение
конкретного прогона к провалу. Для выбора задач это не проблема, но метрику
надо называть правильно.

\textbf{Сигналы, требующие прогонов, нам не подходят.}
Hathidara~\cite{hathidara2026failure} калибрует уверенность агента по исходу
эпизода и получает AUROC $0{,}68$--$0{,}85$. Mehta~\cite{mehta2026disagree}
использует разброс траекторий: 2{,}3--4{,}2 различных последовательности
действий на 10 запусков, AUROC $0{,}62$--$0{,}78$, точность $82$--$87\%$ на
согласованных задачах против $41$--$65\%$ на несогласованных. Оба сигнала
требуют нескольких прогонов, то есть тратят тот самый бюджет, который нужно
экономить.

\subsection{Психометрика и IRT}

IRT — самый распространённый формальный аппарат для «сложности пункта».
Обзоры~\cite{alkhuzaey2024text,peters2025textbased} систематизируют
автоматическое предсказание трудности по тексту: признаки, модели, метрики.
Второй обзор покрывает 37 статей по май 2025 года и отмечает переход от
ручных признаков к трансформерам.

\textbf{Что работает на практике.}
Štěpánek и соавторы~\cite{stepanek2023item} сравнивают алгоритмы на текстовых
признаках: elastic net даёт лучший RMSE, random forest — лучшую точность
классификации; ML конкурирует с экспертами. Li и
соавторы~\cite{li2025item} показывают, что дообученные малые модели
(BERT, RoBERTa) превосходят победителя BEA~2024 по RMSE, а GPT-4 предсказывает
трудность плохо. Zotos и соавторы~\cite{zotos2024doubtful} добавляют признаки
неуверенности нескольких LLM и получают SOTA на USMLE и CMCQRD.

\textbf{Симуляция отвечает лучше, чем прямой вопрос.}
Acquaye и соавторы~\cite{acquaye2026calculators} просят LLM разыграть
«класс» учеников разного уровня и по ответам строят IRT. Корреляции с реальной
трудностью — $0{,}75$, $0{,}76$ и $0{,}82$ для 4, 8 и 12 классов; при этом
математически более слабая Gemma предсказывает лучше сильных Llama и Qwen.
SMART~\cite{scarlatos2025smart} делает то же с DPO-выравниванием симулированных
студентов и превосходит baselines на двух наборах. Прямая оценка сложности
сильной моделью, наоборот, систематически смещена: Li и
соавторы~\cite{li2026struggles} на 20+ моделях видят «машинный консенсус»
вместо человеческой трудности и отсутствие интроспекции собственных границ.

\textbf{IRT на ответах моделей.}
TaskEval~\cite{tambon2026taskeval} применяет IRT к HumanEval+ и ClassEval с
восемью LLM и разными формулировками промпта: сложность задачи зависит от
формулировки, а оценки людей и LLM расходятся. PSN-IRT~\cite{zhou2025lost}
на 11 бенчмарках и 41\,871 пункте показывает, что измерение качества
бенчмарков нестабильно. IRT-Router~\cite{song2025irtrouter} переносит ту же
модель на маршрутизацию запросов (20 LLM, 12 датасетов) и работает в cold
start.

Главный вывод группы: IRT даёт интерпретируемую латентную сложность, но
требует матрицы ответов. Для новой задачи без ответов нужен либо перенос
признаков (как в~\cite{ge2026psychometrics}), либо симуляция (как
в~\cite{acquaye2026calculators,scarlatos2025smart}).

\subsection{Код и задачи Programming Online Judge}

Задачи по программированию — ближайший к нам домен.

Kim и соавторы~\cite{kim2023problem} собрали AMT — крупный набор задач
Codeforces — и предсказывают одновременно тег алгоритма и сложность; это
первая работа именно про сложность алгоритмических задач. Wang и
соавторы~\cite{wang2024estimating} добавляют к тексту пример кода и получают
выигрыш от обеих модальностей на двух наборах POJ. CodeElo~\cite{quan2025codeelo}
строит Elo-рейтинг сложности, сверенный с платформой, и даёт
человекосопоставимую шкалу. TaskEval~\cite{tambon2026taskeval}, рассмотренный
выше, показывает, что даже для кода сложность зависит от формулировки
промпта, а не только от задачи.

Общее ограничение: во всех работах есть либо рейтинг платформы, либо пример
решения. В соревновании пример решения может отсутствовать, а рейтинга
сложности нет.

\subsection{Маршрутизация и каскады}

Маршрутизация отвечает не «решать ли задачу», а «какой моделью её решать».
Это другая задача, но методы оценки сложности запроса переносятся.

RouteLLM~\cite{ong2024routellm} обучает роутер на preference data и в отдельных
случаях снижает стоимость более чем в два раза без потери
качества. FrugalGPT~\cite{chen2023frugalgpt}
строит каскад моделей и экономит до 98\% при качестве GPT-4.
IRT-Router~\cite{song2025irtrouter} даёт интерпретируемую сложность запроса.
Adaptive LLM Routing~\cite{panda2025adaptive} формулирует выбор модели как
contextual bandit с бюджетом и не требует полной матрицы качеств.

Две работы ближе всего к априорной оценке. VeriDispatcher~\cite{wang2025veridispatcher}
обучает компактный классификатор сложности на эмбеддингах описания RTL-задачи
и получает +18\% точности на RTLLM, используя лишь 40\% коммерческих вызовов,
а на VerilogEval снижает расход на 25\%. Классификатор сложности на
BERT-эмбеддингах~\cite{harshith2025cost} даёт weighted F1~$=0{,}89$; вместе с
правилами маршрутизации стоимость падает с \$12{,}08 до \$0{,}34 на 1000
промптов.

Полезные предупреждения. Lai и соавторы~\cite{lai2026overpay} показывают, что
при свободном бюджете роутеры переплачивают, выбирая сильную модель там, где
слабая даёт тот же результат. Garg и Sagtani~\cite{garg2026unsolvability} на
206\,000 пар «запрос--модель» показывают, что значительная часть «нерешаемости»
— артефакт судьи, усечения и формата, а не свойство задачи. Значит, и метка
«невыполнимо» в наших данных требует проверки, а не слепого доверия.

\subsection{Бюджет, порядок и распределение}

Эта группа отвечает на вопрос, что делать с оценкой после того, как она
получена.

TRIAGE~\cite{nazi2026triage} — самая близкая постановка. Модель получает пул
задач и бюджет в токенах и должна до исполнения выбрать задачи, порядок и
долю бюджета на каждую; план сравнивается с oracle по triage efficiency.
Современные модели показывают большой разрыв с oracle на математике, науке и
коде. Fan и соавторы~\cite{fan2026thinking} независимо подтверждают: при
общем бюджете модели ведут себя как жадные последовательные решатели,
чувствительны к порядку и игнорируют баллы задач.

Формальные методы распределения дают выигрыш. ZEBRA~\cite{hamri2026zebra}
сводит разбиение бюджета по фазам пайплайна к нелинейному knapsack и при
половине неограниченного бюджета возвращает 94{,}4\% качества против 88{,}1\%
у прямого распределения LLM. Контроллер на Value-of-Information~\cite{fang2026voi}
выбирает следующее действие под двойным бюджетом (вызовы инструментов и
токены) и даёт устойчивый прирост на четырёх бенчмарках. TALE~\cite{han2025tale}
подбирает бюджет рассуждения по сложности задачи и снижает затраты при
небольшом падении точности. Bandit-подход~\cite{zuo2025bandit} оценивает
сложность на лету и даёт до +11{,}10\% (15{,}04\% относительно) на MATH-500 и
+10{,}82\% на AIME25, причём учится не тратить вычисления на нерешаемые
задачи. Snell и соавторы~\cite{snell2024scaling} показывают, что оптимальное
распределение вычислений зависит от сложности запроса и превосходит
равномерное.

Вывод группы: считать оценку бесполезной без решающего правила нельзя.
Нужны и предиктор, и политика выбора; по отдельности каждая часть даёт мало.

\subsection{Выполнимость и отказ}

Невыполнимые задачи — отдельная проблема.

Peng и соавторы~\cite{peng2026solvability} вводят UnsolvableQA и обучают
разделять «в запросе не хватает данных» и «модель не умеет»: на Qwen3-4B
детекция нерешаемости превышает 85\%, точность на решаемых задачах растёт с
43{,}4\% до 69{,}4\%. The Impossible Test~\cite{noever2024impossible}
показывает, что на 675 принципиально нерешаемых задачах 12 LLM признают
незнание лишь в 62--68\% случаев, причём на более трудных задачах — чаще.
Kadavath и соавторы~\cite{kadavath2022know} показывают, что модели неплохо
калиброваны на знакомых форматах (P(True), P(IK)), но плохо переносят
калибровку на новые задачи.

Для харнесса это означает: детектор невыполнимости нужен отдельный и
обучаемый, а не порог по уверенности модели.

\subsection{Классическая аналогия: время работы алгоритма}

Hutter и соавторы~\cite{hutter2014runtime} систематизируют предсказание
времени работы алгоритма на экземпляре: признаки экземпляра, регрессия,
распределения времени, censored data. Эта постановка старше LLM, но
структурно совпадает с нашей: есть экземпляр (условие задачи), есть
стохастическая стоимость с тяжёлым хвостом, есть выбор алгоритма (модели).
Оттуда переносятся три практики: предсказывать распределение, а не только
среднее; учитывать цензурирование (задача не досчитана); оценивать не MSE, а
выгоду решения.

\section{Сравнение}

Таблица~\ref{tab:sources} сводит включённые работы по цели, подходу, данным и
ограничению. Строки сгруппированы по типу цели: успех ($p$), стоимость
($C$), выполнимость ($F$), общая сложность.

{\footnotesize
\setlength{\tabcolsep}{3pt}
\renewcommand{\arraystretch}{1.15}
\begin{longtable}{@{}p{0.7cm}p{3.0cm}p{2.2cm}p{3.1cm}p{3.5cm}p{2.3cm}@{}}
\caption{Включённые работы.}\label{tab:sources}\\
\toprule
ID & Работа & Цель & Подход & Данные и результат & Ограничение \\
\midrule
\endfirsthead
\toprule
ID & Работа & Цель & Подход & Данные и результат & Ограничение \\
\midrule
\endhead
\bottomrule
\endfoot

S01 & Krsteski, Meyer \cite{krsteski2026difficulty} & сложность & предикторы без прогонов, энтропия & 17 бенчмарков; AUC маскирует ошибки & препринт \\
S02 & Al-Haque, Johnson \cite{alhaque2026difficult} & $p$ & статические признаки патча и репозитория & CoderForge; AUC 0{,}863 & один датасет \\
S03 & Ge et al. \cite{ge2026psychometrics} & $p$ & IRT + признаки задачи, LLM/scaffold & невиданные задачи и комбинации & нужна история оценок \\
S04 & Yin et al. \cite{yin2025tasksense} & сложность & когнитивная цепочка + регрессия & GUI; $R^2=0{,}46$ & GUI-домен \\
S05 & Hong et al. \cite{hong2026circe} & $p$ & граф ограничений + GNN & AUROC 0{,}691$\to$0{,}727 & нужны явные ограничения \\
S06 & Lutati \cite{lutati2025capability} & $C$ & информационные границы & $C_{\text{eff}}$ нижне ограничивает затраты & нужны $I_{\text{total}}$, $I_{\text{step}}$ \\
S07 & EsfandyariDoulabi et al. \cite{esfandiaridoulabi2026disentangling} & метрика & анализ пар прогонов & $>$99{,}93\% пар межзадачные & не предиктор \\
S08 & Hathidara \cite{hathidara2026failure} & $p$ & калибровка уверенности & AUROC 0{,}68--0{,}85 & нужен RL \\
S09 & Mehta \cite{mehta2026disagree} & $p$ & разброс траекторий & AUROC 0{,}62--0{,}78 & нужны повторные прогоны \\
S10 & AlKhuzaey et al. \cite{alkhuzaey2024text} & обзор & систематический обзор & признаки и модели & образование \\
S11 & Peters et al. \cite{peters2025textbased} & обзор & 37 статей & переход к трансформерам & крупные оценивания \\
S12 & Tambon et al. \cite{tambon2026taskeval} & сложность & IRT + варианты промпта & HumanEval+, ClassEval & нужны прогоны моделей \\
S13 & Zhou et al. \cite{zhou2025lost} & параметры & PSN-IRT & 41\,871 пункт & нужна матрица ответов \\
S14 & Acquaye et al. \cite{acquaye2026calculators} & сложность & симуляция учеников + IRT & $r=0{,}75$--$0{,}82$ & школьная математика \\
S15 & Li et al. \cite{li2026struggles} & сложность & 20+ моделей, симуляция уровня & рассогласование, нет интроспекции & отрицательный результат \\
S16 & Li et al. \cite{li2025item} & сложность & fine-tuned BERT, GPT-4 & малые модели лучше GPT-4 & нужна разметка \\
S17 & Zotos et al. \cite{zotos2024doubtful} & сложность & неуверенность + текст, RF & SOTA на USMLE, CMCQRD & MCQ \\
S18 & Štěpánek et al. \cite{stepanek2023item} & сложность & elastic net, RF, SVM & RMSE, конкуренция с экспертами & ручные признаки \\
S19 & Kim et al. \cite{kim2023problem} & сложность & многозадачная модель & AMT, Codeforces & рейтинг как метка \\
S20 & Wang et al. \cite{wang2024estimating} & сложность & текст + код & два набора POJ & нужен пример решения \\
S21 & Quan et al. \cite{quan2025codeelo} & шкала & Elo & Codeforces & узкий домен \\
S22 & Song et al. \cite{song2025irtrouter} & $p$ & IRT + warm-up & 20 LLM, 12 датасетов & нет стоимости \\
S23 & Ong et al. \cite{ong2024routellm} & выбор модели & preference-роутер & $>$2$\times$ дешевле & не выбор задачи \\
S24 & Chen et al. \cite{chen2023frugalgpt} & каскад & LLM cascade & до 98\% экономии & по ответам \\
S25 & Panda et al. \cite{panda2025adaptive} & выбор модели & contextual bandit & адаптация при бюджете & бюджет на запрос \\
S26 & Wang et al. \cite{wang2025veridispatcher} & сложность & эмбеддинг + классификатор & +18\% при 40\% вызовов & RTL \\
S27 & Harshith et al. \cite{harshith2025cost} & сложность & BERT + LightGBM & F1 0{,}89; $-\,$95\% цены & синтетика \\
S28 & Lai et al. \cite{lai2026overpay} & выбор модели & cost-aware ranking & переплата при свободном бюджете & не оценка задачи \\
S29 & Garg, Sagtani \cite{garg2026unsolvability} & $F$ & декомпозиция артефактов & 206\,000 пар & не предиктор \\
S30 & Nazi, Dipta \cite{nazi2026triage} & план & prospective control & разрыв с oracle & оценка способности \\
S31 & Fan et al. \cite{fan2026thinking} & план & экзаменационный протокол & жадность, игнор баллов & нет метода \\
S32 & Hamri, Talgam-Cohen \cite{hamri2026zebra} & бюджет & knapsack + water-filling & 94{,}4\% против 88{,}1\% & по фазам \\
S33 & Fang et al. \cite{fang2026voi} & бюджет & Value-of-Information & 4 бенчмарка & онлайн-сигнал \\
S34 & Han et al. \cite{han2025tale} & $C$ & бюджет рассуждения & ниже затраты, небольшое падение & одна задача \\
S35 & Zuo, Zhu \cite{zuo2025bandit} & бюджет & bandit & +11{,}10\% на MATH-500 & нужен фидбек \\
S36 & Snell et al. \cite{snell2024scaling} & бюджет & compute-optimal scaling & зависит от сложности & одношаговые задачи \\
S37 & Peng et al. \cite{peng2026solvability} & $F$ & UnsolvableQA + RL & $>$85\% детекции & свой бенчмарк \\
S38 & Noever, McKee \cite{noever2024impossible} & $F$ & опрос 12 LLM & 62--68\% признания & только признание \\
S39 & Kadavath et al. \cite{kadavath2022know} & $p$ & P(True), P(IK) & калибровка на знакомом & плохо на новом \\
S40 & Hutter et al. \cite{hutter2014runtime} & $C$ & runtime prediction & распределения, цензурирование & не LLM \\
S41 & Scarlatos et al. \cite{scarlatos2025smart} & сложность & DPO + IRT & лучше baselines & нужна калибровка \\
\end{longtable}
}

\section{Пробел и вывод}

\subsection{Что не покрыто}

\begin{enumerate}
  \item \textbf{Совместная оценка $p$ и $C$.} Работы предсказывают либо
        вероятность успеха, либо стоимость. Единственная априорная оценка
        стоимости~\cite{lutati2025capability} теоретическая и требует
        информационной структуры задачи. Для выбора задач нужна именно пара
        «ожидаемый балл / ожидаемые токены».
  \item \textbf{Агентский код при общем бюджете.} IRT и психометрика работают
        с матрицей ответов; маршрутизация выбирает модель на один запрос;
        распределение бюджета изучается для фаз пайплайна или одного
        запроса. Постановка «100 задач, один бюджет, скрытые тесты,
        15--25\% невыполнимых» в литературе отдельно не измерялась.
  \item \textbf{Cold start без истории.} Почти все сильные результаты
        опираются на прошлые прогоны или рейтинги платформы. Для новых задач
        остаются признаки условия и перенос с других бенчмарков, а
        переносимость между моделями ограничена~\cite{ge2026psychometrics,li2026struggles}.
  \item \textbf{Метрика качества оценки.} AUC может маскировать плохие
        оценки~\cite{krsteski2026difficulty}, а pooled AUROC завышен из-за
        межзадачных пар~\cite{esfandiaridoulabi2026disentangling}. Общепринятой
        метрики «качество оценки для бюджетного выбора» нет.
  \item \textbf{Выполнимость отдельно от сложности.} Детекторы
        нерешаемости~\cite{peng2026solvability,noever2024impossible} не
        оценивают стоимость, а «нерешаемость» в бенчмарках частично является
        артефактом оценки~\cite{garg2026unsolvability}.
\end{enumerate}

\subsection{Что из этого следует для решения}

Литература не даёт готового предиктора, но даёт порядок действий.

\begin{enumerate}
  \item \textbf{Определить цель.} Оценивать две величины: распределение
        стоимости (логарифм токенов, квантили) и вероятность успеха; отдельно
        — бинарную выполнимость. Сложность для выбора — это не число, а
        ожидаемая выгода на токен.
  \item \textbf{Брать признаки, доступные до исполнения.} Длина и структура
        условия, число и типы файлов, наличие тестов, размер и связность
        репозитория, язык, близость к известным задачам. Не использовать
        решение, тесты и скрытые состояния: в финале их нет.
  \item \textbf{Не спрашивать сложность у сильной LLM напрямую.} Прямая
        оценка смещена~\cite{li2026struggles,li2025item}. LLM полезна как
        аннотатор признаков~\cite{peters2025textbased} или как симулятор
        трудности~\cite{acquaye2026calculators}, но не как оракул.
  \item \textbf{Учить дешёвую модель на своих траекториях.} Малые
        дообученные модели и градиентный бустинг на признаках дают хорошее
        отношение качества к цене~\cite{li2025item,stepanek2023item,harshith2025cost}.
        Если моделей несколько, IRT-разложение помогает отделить способность
        модели от scaffold~\cite{ge2026psychometrics}.
  \item \textbf{Калибровать и проверять не только AUC.} Нужны ранговая
        корреляция, MAE на логарифме стоимости, калибровка вероятностей и
        итоговая выгода на фиксированном бюджете. Оценивать на отложенных
        задачах, а не на pooled-прогонах.
  \item \textbf{Добавить политику выбора.} Ранжировать задачи по
        ожидаемой выгоде на токен с верхней оценкой (UCB), пропускать задачи
        с низкой вероятностью выполнимости, обновлять оценки после каждой
        решённой задачи. Bandit и knapsack дают измеримый выигрыш в близких
        постановках~\cite{panda2025adaptive,hamri2026zebra,zuo2025bandit}.
  \item \textbf{Ограничить стоимость оценки.} Оценка должна стоить доли
        процента бюджета: статические признаки плюс один дешёвый вызов.
        Сигналы, требующие повторных прогонов~\cite{mehta2026disagree},
        для этого слишком дороги.
\end{enumerate}

Проверять это надо экспериментом, а не сразу спецификацией: сначала собрать
траектории на доступных задачах, обучить предиктор, сравнить три стратегии
(текущий порядок, эвристика «балл / длина условия», обученный предиктор) при
одинаковом бюджете. Запись эксперимента создаётся в
\texttt{research/experiments/} до запуска. Если прирост выгоды на токен
устойчив, писать спецификацию изменения
\texttt{04\_budget\_and\_selection.md}.

\subsection{Короткий ответ}

По условию задачи сложность оценивается не одним числом, а тремя
предсказаниями: вероятность успеха, распределение стоимости в токенах и
вероятность выполнимости. Лучше всего работают структурные признаки условия —
размер и фрагментация предполагаемого изменения, масштаб репозитория, число
файлов и тестов, — а не длина текста и не прямая оценка сильной LLM. Модель
надо обучать на собственных траекториях и калибровать на отложенных задачах;
готового переносимого предиктора литература не даёт. Решение о задаче
принимается не по сложности, а по ожидаемой выгоде на токен с учётом
неопределённости; невыполнимость предсказывается отдельным классификатором.
Стоимость самой оценки должна быть пренебрежимо мала по сравнению с бюджетом.

\section{Ограничения обзора}

\begin{itemize}
  \item Две базы вместо четырёх: Semantic Scholar и Google Scholar вернули
        HTTP 403. Возможны пропуски, особенно в индустриальных публикациях.
  \item Просматривались верхние 8--25 записей на запрос, а не вся выдача.
  \item Для большинства работ данные извлечены из аннотаций; полные тексты
        читались выборочно. Числа в таблице — из аннотаций и могут отличаться
        от итоговых таблиц статей.
  \item Препринты 2026 года не рецензированы и могут измениться.
  \item Метрики несопоставимы: корреляция, AUC, F1, RMSE, экономия денег или
        токенов. Мета-анализ невозможен.
  \item Обзор не проверяет воспроизводимость включённых работ и не оценивает
        риск систематической ошибки публикации.
\end{itemize}

\bibliographystyle{plain}
\bibliography{refs}

\end{document}
